\import{relation/well-founded.tex}
\import{order/zorn.tex}
\import{order/well-order.tex}
\import{function.tex}

\section{Well-founded recursion}

The recursion operator in this section is deliberately bounded.  A step
function has type
$A\times\pow{A\times B}\longrightarrow B$: its second argument is the graph
of the already defined values below the current point.  This formulation is
set-sized, first-order, and sufficient whenever a codomain bound is known.

\begin{definition}\label{recursive_partials}
    $\recursivePartials{A}{B}{R}{H} =
        \{f\in\pow{A\times B}\mid
            \text{$f\in\funs{\dom{f}}{B}$ and
            $\dom{f}$ is an initial segment of $A$ with respect to $R$ and
            for every $x\in\dom{f}$ we have
            $f(x)=H((x,\restrl{f}{\downward{R}{x}}))$}\}$.
\end{definition}

\begin{definition}\label{recursive_functions}
    $\recursiveFunctions{A}{B}{R}{H}=
        \{f\in\funs{A}{B}\mid
            \text{for every $x\in A$ we have
            $f(x)=H((x,\restrl{f}{\downward{R}{x}}))$}\}$.
\end{definition}

\begin{proposition}\label{empty_recursive_partial}
    $\emptyset\in\recursivePartials{A}{B}{R}{H}$.
\end{proposition}
\begin{proof}
    $\emptyset\in\pow{A\times B}$ by \obv{powerset_bottom}.
    $\emptyset\in\funs{\emptyset}{B}$
        by \obv{emptyset_is_function_on_emptyset,codom_of_emptyset_can_be_anything,funs_intro}.
    $\dom{\emptyset}=\emptyset$ by \obv{dom_emptyset}.
    $\emptyset$ is an initial segment of $A$ with respect to $R$
        by \obv{emptyset_initial_segment}.
    Follows by \auto{recursive_partials,emptyset}.
\end{proof}

\begin{proposition}\label{recursive_partials_inhabited}
    $\recursivePartials{A}{B}{R}{H}$ is inhabited.
\end{proposition}
\begin{proof}
    Follows by \auto{empty_recursive_partial}.
\end{proof}

\begin{proposition}\label{recursive_partial_is_function}
    Let $f\in\recursivePartials{A}{B}{R}{H}$.
    Then $f$ is a function and $f\subseteq A\times B$.
\end{proposition}
\begin{proof}
    $f\in\pow{A\times B}$ and $f\in\funs{\dom{f}}{B}$
        by \obv{recursive_partials}.
    Follows by \obv{powerset_elim,funs_is_function}.
\end{proof}

\begin{proposition}\label{recursive_partial_domain_initial}
    Let $f\in\recursivePartials{A}{B}{R}{H}$.
    Then $\dom{f}$ is an initial segment of $A$ with respect to $R$.
\end{proposition}
\begin{proof}
    Follows by \obv{recursive_partials}.
\end{proof}

\begin{proposition}\label{recursive_partial_equation}
    Let $f\in\recursivePartials{A}{B}{R}{H}$.
    Let $x\in\dom{f}$.
    Then $f(x)=H((x,\restrl{f}{\downward{R}{x}}))$.
\end{proposition}
\begin{proof}
    Follows by \obv{recursive_partials}.
\end{proof}

\begin{proposition}\label{recursive_partial_union_is_function}
    Suppose $C\subseteq\recursivePartials{A}{B}{R}{H}$.
    Suppose $C$ is an inclusion chain.
    Then $\unions{C}$ is a function and $\unions{C}\subseteq A\times B$.
\end{proposition}
\begin{proof}
    Every element of $C$ is a function
        by \auto{subseteq,recursive_partial_is_function}.
    For all $f,g\in C$ we have $f\subseteq g$ or $g\subseteq f$
        by \obv{inclusion_chain}.
    Thus $\unions{C}$ is a function
        by \obv{unions_of_compatible_family_of_function_is_function}.
    $\unions{C}\subseteq A\times B$
        by \auto{subseteq,unions_iff,recursive_partial_is_function}.
\end{proof}

\begin{proposition}\label{recursive_partial_union_domain_initial}
    Suppose $C\subseteq\recursivePartials{A}{B}{R}{H}$.
    Then $\dom{\unions{C}}$ is an initial segment of $A$ with respect to $R$.
\end{proposition}
\begin{proof}
    $\dom{\unions{C}}\subseteq A$
        by \auto{subseteq,dom_iff,unions_iff,recursive_partial_is_function,times_tuple_elim}.
    For all $x\in\dom{\unions{C}}$ we have for all $y\in A$
        such that $y\mathrel{R}x$ we have $y\in\dom{\unions{C}}$
        by \auto{dom_iff,unions_iff,subseteq,recursive_partial_domain_initial,initial_segment}.
    Follows by \obv{initial_segment}.
\end{proof}

\begin{lemma}\label{recursive_partial_union_restriction}
    Suppose $C\subseteq\recursivePartials{A}{B}{R}{H}$.
    Suppose $C$ is an inclusion chain.
    Suppose $R$ is a binary relation on $A$.
    Let $f\in C$.
    Let $x\in\dom{f}$.
    Then
        $\restrl{\unions{C}}{\downward{R}{x}}
        =\restrl{f}{\downward{R}{x}}$.
\end{lemma}
\begin{proof}
    Let $g=\unions{C}$.
    $f$ and $g$ are functions
        by \auto{subseteq,recursive_partial_is_function,recursive_partial_union_is_function}.
    $f\subseteq g$ by \auto{subseteq,unions_iff}.
    Show that for every $y\in\downward{R}{x}$ such that
        $y\in\dom{f}\union\dom{g}$ we have
        $y\in\dom{f}\inter\dom{g}$ and $f(y)=g(y)$.
    \begin{subproof}
        Fix $y$.
        Assume $y\in\downward{R}{x}$ and $y\in\dom{f}\union\dom{g}$.
        $y\mathrel{R}x$ by \obv{downward_closure_iff}.
        $y\in A$ by \obv{subseteq,times_tuple_elim}.
        $\dom{f}$ is an initial segment of $A$ with respect to $R$
            by \obv{subseteq,recursive_partial_domain_initial}.
        Hence $y\in\dom{f}$ by \obv{initial_segment}.
        Thus $y\in\dom{g}$ by \auto{subseteq,dom_iff}.
        Hence $y\in\dom{f}\inter\dom{g}$ by \obv{inter_intro}.
        Follows by \obv{subfunction_apply}.
    \end{subproof}
    Follows by \auto{restrictions_agree}.
\end{proof}

\begin{proposition}\label{recursive_partial_chain_union}
    Suppose $C\subseteq\recursivePartials{A}{B}{R}{H}$.
    Suppose $C$ is an inclusion chain.
    Suppose $R$ is a binary relation on $A$.
    Then $\unions{C}\in\recursivePartials{A}{B}{R}{H}$.
\end{proposition}
\begin{proof}
    Let $g=\unions{C}$.
    $g$ is a function and $g\subseteq A\times B$
        by \auto{recursive_partial_union_is_function}.
    Thus $g\in\pow{A\times B}$ by \obv{powerset_intro}.
    For every $x\in\dom{g}$ we have $g(x)\in B$
        by \auto{function_apply_elim,subseteq,times_tuple_elim}.
    Thus $g$ is a function to $B$.
    Hence $g\in\funs{\dom{g}}{B}$ by \obv{funs_intro}.
    $\dom{g}$ is an initial segment of $A$ with respect to $R$
        by \auto{recursive_partial_union_domain_initial}.
    Show that for every $x\in\dom{g}$ we have
        $g(x)=H((x,\restrl{g}{\downward{R}{x}}))$.
    \begin{subproof}
        Fix $x\in\dom{g}$.
        Take $f\in C$ such that $x\in\dom{f}$
            by \auto{dom_iff,unions_iff,relation,function_member_elim,pair_eq_iff}.
        $f\in\recursivePartials{A}{B}{R}{H}$ by \obv{subseteq}.
        $f\subseteq g$ by \auto{subseteq,unions_iff}.
        $f(x)=g(x)$ by \obv{subfunction_apply,recursive_partial_is_function}.
        $f(x)=H((x,\restrl{f}{\downward{R}{x}}))$
            by \obv{recursive_partial_equation}.
        $\restrl{g}{\downward{R}{x}}=\restrl{f}{\downward{R}{x}}$
            by \auto{recursive_partial_union_restriction}.
    \end{subproof}
    Follows by \auto{recursive_partials}.
\end{proof}

\begin{theorem}\label{maximal_recursive_partial}
    Suppose $R$ is a binary relation on $A$.
    Then there exists $M\in\recursivePartials{A}{B}{R}{H}$ such that
        for every $f\in\recursivePartials{A}{B}{R}{H}$ such that
            $M\subseteq f$ we have $f=M$.
\end{theorem}
\begin{proof}
    Let $P=\recursivePartials{A}{B}{R}{H}$.
    $P$ is inhabited by \auto{recursive_partials_inhabited}.
    Show that for every $C$ such that $C\subseteq P$ and $C$ is inhabited
        and $C$ is an inclusion chain we have $\unions{C}\in P$.
    \begin{subproof}
        Fix $C$.
        Assume $C\subseteq P$ and $C$ is inhabited and $C$ is an inclusion chain.
        Follows by \obv{recursive_partial_chain_union}.
    \end{subproof}
    Follows by \auto{zorn_subset_family}.
\end{proof}

\begin{lemma}\label{recursive_extension_restriction}
    Suppose $R$ is a binary relation on $A$.
    Let $f\in\recursivePartials{A}{B}{R}{H}$.
    Let $x\in A\setminus\dom{f}$.
    Suppose for every $y\in A$ such that $y\mathrel{R}x$ we have $y\in\dom{f}$.
    Let $b\in B$.
    Let $g=\cons{(x,b)}{f}$.
    Let $z\in\dom{g}$.
    Then $\restrl{g}{\downward{R}{z}}=\restrl{f}{\downward{R}{z}}$.
\end{lemma}
\begin{proof}
    $f$ is a function by \auto{recursive_partial_is_function}.
    $x\notin\dom{f}$ by \obv{setminus}.
    $g\in\funs{\cons{x}{\dom{f}}}{B}$
        by \auto{recursive_partials,funs,extend_function_at_new_point}.
    Hence $g$ is a function and $\dom{g}=\cons{x}{\dom{f}}$
        by \obv{funs_is_function,funs}.
    $f\subseteq g$ by \auto{subseteq,cons_iff}.
    Show that for every $y\in\downward{R}{z}$ such that
        $y\in\dom{g}\union\dom{f}$ we have
        $y\in\dom{g}\inter\dom{f}$ and $g(y)=f(y)$.
    \begin{subproof}
        Fix $y$.
        Assume $y\in\downward{R}{z}$ and $y\in\dom{g}\union\dom{f}$.
        $y\mathrel{R}z$ by \obv{downward_closure_iff}.
        $y\in A$ by \obv{subseteq,times_tuple_elim}.
        $z=x$ or $z\in\dom{f}$ by \auto{cons_iff}.
        $\dom{f}$ is an initial segment of $A$ with respect to $R$
            by \obv{recursive_partial_domain_initial}.
        Hence $y\in\dom{f}$ by \auto{initial_segment}.
        Thus $y\in\dom{g}$ by \auto{subseteq,dom_iff}.
        Hence $y\in\dom{g}\inter\dom{f}$ by \obv{inter_intro}.
        Follows by \obv{subfunction_apply}.
    \end{subproof}
    Follows by \obv{restrictions_agree}.
\end{proof}

\begin{proposition}\label{extend_recursive_partial}
    Suppose $R$ is a binary relation on $A$.
    Let $H\in\funs{A\times\pow{A\times B}}{B}$.
    Let $f\in\recursivePartials{A}{B}{R}{H}$.
    Let $x\in A\setminus\dom{f}$.
    Suppose for every $y\in A$ such that $y\mathrel{R}x$ we have $y\in\dom{f}$.
    Let $p=\restrl{f}{\downward{R}{x}}$.
    Let $b=H((x,p))$.
    Then $\cons{(x,b)}{f}\in\recursivePartials{A}{B}{R}{H}$.
\end{proposition}
\begin{proof}
    $f$ is a function and $f\subseteq A\times B$
        by \obv{recursive_partial_is_function}.
    $p\subseteq f$ by \auto{restrl_subseteq}.
    Hence $p\subseteq A\times B$ by \obv{subseteq_transitive}.
    Thus $p\in\pow{A\times B}$ by \obv{powerset_intro}.
    Hence $(x,p)\in A\times\pow{A\times B}$ by \obv{setminus,times_tuple_intro}.
    Thus $b\in B$ by notion.
    Let $g=\cons{(x,b)}{f}$.
    $x\notin\dom{f}$ by \obv{setminus}.
    $f\in\funs{\dom{f}}{B}$ by \obv{recursive_partials}.
    Hence $g\in\funs{\cons{x}{\dom{f}}}{B}$
        by \auto{extend_function_at_new_point}.
    Thus $g$ is a function and $\dom{g}=\cons{x}{\dom{f}}$
        by \obv{funs_is_function,funs}.
    $(x,b)\in A\times B$ by \obv{setminus,times_tuple_intro}.
    $g\subseteq A\times B$ by \auto{cons_iff,subseteq}.
    Thus $g\in\pow{A\times B}$ by \obv{powerset_intro}.
    Show that $\dom{g}$ is an initial segment of $A$ with respect to $R$.
    \begin{subproof}
        $\dom{g}\subseteq A$
            by \auto{cons_iff,setminus,recursive_partial_domain_initial,initial_segment,subseteq}.
        Show that for all $z,y$ such that $z\in\dom{g}$ and $y\in A$
            and $y\mathrel{R}z$ we have $y\in\dom{g}$.
        \begin{subproof}
            Fix $z,y$.
            Assume $z\in\dom{g}$ and $y\in A$ and $y\mathrel{R}z$.
            $z=x$ or $z\in\dom{f}$ by \auto{cons_iff}.
            $\dom{f}$ is an initial segment of $A$ with respect to $R$
                by \obv{recursive_partial_domain_initial}.
            Thus $y\in\dom{f}$ by \auto{initial_segment}.
            Follows by \auto{cons_iff}.
        \end{subproof}
        Follows by \obv{initial_segment}.
    \end{subproof}
    Show that for every $z\in\dom{g}$ we have
        $g(z)=H((z,\restrl{g}{\downward{R}{z}}))$.
    \begin{subproof}
        Fix $z\in\dom{g}$.
        $z=x$ or $z\in\dom{f}$ by \auto{cons_iff}.
        $\restrl{g}{\downward{R}{z}}=\restrl{f}{\downward{R}{z}}$
            by \obv{recursive_extension_restriction}.
        \begin{byCase}
            \caseOf{$z=x$.}
                $(x,b)\in g$ by \obv{cons_left}.
                Hence $g(z)=b$ by \obv{function_apply_intro}.
                $\restrl{g}{\downward{R}{z}}=p$.
            \caseOf{$z\neq x$.}
                $z\in\dom{f}$.
                $f\subseteq g$ by \auto{subseteq,cons_iff}.
                Hence $f(z)=g(z)$ by \obv{subfunction_apply}.
                $f(z)=H((z,\restrl{f}{\downward{R}{z}}))$
                    by \obv{recursive_partial_equation}.
        \end{byCase}
    \end{subproof}
    Follows by \auto{recursive_partials}.
\end{proof}

\begin{theorem}\label{bounded_well_founded_recursion_exists}
    Suppose $R$ is a binary relation on $A$ and $R$ is well-founded on $A$.
    Let $H\in\funs{A\times\pow{A\times B}}{B}$.
    Then $\recursiveFunctions{A}{B}{R}{H}$ is inhabited.
\end{theorem}
\begin{proof}
    Take $M\in\recursivePartials{A}{B}{R}{H}$ such that
        for every $f\in\recursivePartials{A}{B}{R}{H}$ such that
            $M\subseteq f$ we have $f=M$
        by \obv{maximal_recursive_partial}.
    $\dom{M}$ is an initial segment of $A$ with respect to $R$
        by \obv{recursive_partial_domain_initial}.
    Show that $\dom{M}=A$.
    \begin{subproof}
        Suppose not.
        Let $D=A\setminus\dom{M}$.
        $\dom{M}\subseteq A$ by \obv{initial_segment}.
        Hence $D$ is inhabited
            by \auto{subset,difference_with_proper_subset_is_inhabited}.
        $D\subseteq A$ by \auto{setminus_subseteq}.
        Take $x\in D$ such that for every $y\in D$ we have $y\not\mathrel{R}x$
            by \obv{well_founded_on}.
        For every $y\in A$ such that $y\mathrel{R}x$ we have $y\in\dom{M}$
            by \auto{setminus}.
        Let $p=\restrl{M}{\downward{R}{x}}$.
        Let $b=H((x,p))$.
        Let $g=\cons{(x,b)}{M}$.
        $g\in\recursivePartials{A}{B}{R}{H}$
            by \auto{extend_recursive_partial}.
        $M\subseteq g$ by \auto{subseteq,cons_iff}.
        Hence $g=M$.
        $x\in\dom{g}$ by \auto{dom_cons,cons_left}.
        $x\notin\dom{M}$ by \auto{setminus}.
        Contradiction.
    \end{subproof}
    $M\in\funs{\dom{M}}{B}$ by \obv{recursive_partials}.
    Hence $M\in\funs{A}{B}$.
    For every $x\in A$ we have
        $M(x)=H((x,\restrl{M}{\downward{R}{x}}))$
        by \obv{recursive_partial_equation}.
    Thus $M\in\recursiveFunctions{A}{B}{R}{H}$
        by \obv{recursive_functions}.
\end{proof}

\begin{theorem}\label{bounded_well_founded_recursion_unique}
    Suppose $R$ is a binary relation on $A$ and $R$ is well-founded on $A$.
    Let $H\in\funs{A\times\pow{A\times B}}{B}$.
    Let $f,g\in\recursiveFunctions{A}{B}{R}{H}$.
    Then $f=g$.
\end{theorem}
\begin{proof}
    $f,g\in\funs{A}{B}$ by \obv{recursive_functions}.
    Hence $f$ and $g$ are functions by \obv{funs_is_function}.
    $\dom{f}=A$ and $\dom{g}=A$ by \obv{funs}.
    Let $E=\{x\in A\mid f(x)=g(x)\}$.
    Show that $E$ is progressive for $R$ on $A$.
    \begin{subproof}
        Show that for every $x\in A$ we have
            if for every $y\in A$ such that $y\mathrel{R}x$ we have $y\in E$,
            then $x\in E$.
        \begin{subproof}
            Fix $x\in A$.
            Suppose for every $y\in A$ such that $y\mathrel{R}x$ we have $y\in E$.
            Show that
                $\restrl{f}{\downward{R}{x}}=\restrl{g}{\downward{R}{x}}$.
            \begin{subproof}
                Show that for every $y\in\downward{R}{x}$ such that
                    $y\in\dom{f}\union\dom{g}$ we have
                    $y\in\dom{f}\inter\dom{g}$ and $f(y)=g(y)$.
                \begin{subproof}
                    Fix $y$.
                    Assume $y\in\downward{R}{x}$ and
                        $y\in\dom{f}\union\dom{g}$.
                    $y\mathrel{R}x$ by \obv{downward_closure_iff}.
                    Hence $y\in A$ by \obv{subseteq,times_tuple_elim}.
                    Thus $y\in E$.
                    $f(y)=g(y)$.
                    $y\in\dom{f}\inter\dom{g}$ by \obv{inter_intro}.
                \end{subproof}
                Follows by \obv{restrictions_agree}.
            \end{subproof}
            $f(x)=H((x,\restrl{f}{\downward{R}{x}}))$ and
                $g(x)=H((x,\restrl{g}{\downward{R}{x}}))$
                by \obv{recursive_functions}.
            Hence $f(x)=g(x)$.
        \end{subproof}
        Follows by \obv{progressive_on}.
    \end{subproof}
    Hence $A\subseteq E$ by \obv{well_founded_induction_set}.
    For every $x\in A$ we have $f(x)=g(x)$ by \auto{subseteq}.
    Show that $f\subseteq g$.
    \begin{subproof}
        Show that for every $w\in f$ we have $w\in g$.
        \begin{subproof}
            Fix $w\in f$.
            Take $x\in\dom{f}$ such that $w=(x,f(x))$
                by \obv{function_member_elim}.
            Hence $x\in A$.
            Thus $f(x)=g(x)$.
            $(x,g(x))\in g$ by \obv{function_apply_elim}.
        \end{subproof}
        Follows by \obv{subseteq}.
    \end{subproof}
    Show that $g\subseteq f$.
    \begin{subproof}
        Show that for every $w\in g$ we have $w\in f$.
        \begin{subproof}
            Fix $w\in g$.
            Take $x\in\dom{g}$ such that $w=(x,g(x))$
                by \obv{function_member_elim}.
            Hence $x\in A$.
            Thus $f(x)=g(x)$.
            $(x,f(x))\in f$ by \obv{function_apply_elim}.
        \end{subproof}
        Follows by \obv{subseteq}.
    \end{subproof}
    Follows by \obv{subseteq_antisymmetric}.
\end{proof}

\begin{definition}\label{wfrec}
    $\wfrec{A}{B}{R}{H}=\choice{\recursiveFunctions{A}{B}{R}{H}}$.
\end{definition}

This is a total set-valued term selected from the set of bounded recursive
functions.  Its type, equation, and uniqueness laws apply only under the
relation, well-foundedness, and step-map hypotheses stated below.  When those
hypotheses do not make the candidate set inhabited, global choice supplies no
membership law and no canonical off-domain value is claimed.

\begin{theorem}\label{wfrec_is_recursive}
    Suppose $R$ is a binary relation on $A$ and $R$ is well-founded on $A$.
    Let $H\in\funs{A\times\pow{A\times B}}{B}$.
    Then $\wfrec{A}{B}{R}{H}\in\recursiveFunctions{A}{B}{R}{H}$.
\end{theorem}
\begin{proof}
    $\recursiveFunctions{A}{B}{R}{H}$ is inhabited
        by \obv{bounded_well_founded_recursion_exists}.
    Follows by \auto{wfrec,choice_member}.
\end{proof}

\begin{theorem}\label{wfrec_type}
    Suppose $R$ is a binary relation on $A$ and $R$ is well-founded on $A$.
    Let $H\in\funs{A\times\pow{A\times B}}{B}$.
    Then $\wfrec{A}{B}{R}{H}\in\funs{A}{B}$.
\end{theorem}
\begin{proof}
    Follows by \auto{wfrec_is_recursive,recursive_functions}.
\end{proof}

\begin{theorem}\label{wfrec_equation}
    Suppose $R$ is a binary relation on $A$ and $R$ is well-founded on $A$.
    Let $H\in\funs{A\times\pow{A\times B}}{B}$.
    Let $x\in A$.
    Then
        $\apply{\wfrec{A}{B}{R}{H}}{x}=
            H((x,\restrl{\wfrec{A}{B}{R}{H}}{\downward{R}{x}}))$.
\end{theorem}
\begin{proof}
    Follows by \auto{wfrec_is_recursive,recursive_functions}.
\end{proof}

\begin{theorem}\label{wfrec_unique}
    Suppose $R$ is a binary relation on $A$ and $R$ is well-founded on $A$.
    Let $H\in\funs{A\times\pow{A\times B}}{B}$.
    Let $f\in\recursiveFunctions{A}{B}{R}{H}$.
    Then $f=\wfrec{A}{B}{R}{H}$.
\end{theorem}
\begin{proof}
    $\wfrec{A}{B}{R}{H}\in\recursiveFunctions{A}{B}{R}{H}$
        by \obv{wfrec_is_recursive}.
    Follows by \obv{bounded_well_founded_recursion_unique}.
\end{proof}
